% !TeX program = xelatex
\documentclass[UTF8,zihao=-4,a4paper]{ctexart}
\usepackage[a4paper,margin=25mm]{geometry}
\usepackage{amsmath,amssymb,amsthm,mathtools,booktabs,tabularx}
\usepackage{xcolor,hyperref}
\definecolor{accent}{HTML}{153E57}
\hypersetup{colorlinks,linkcolor=accent,urlcolor=accent,pdftitle={谱梯度法的端点动力学与有限记忆收敛}}
\ctexset{section={format=\Large\bfseries\color{accent}},subsection={format=\large\bfseries}}
\setlength{\parindent}{0pt}
\setlength{\parskip}{5pt plus 1pt}
\newtheorem{proposition}{命题}
\renewcommand{\proofname}{证明}
\newcommand{\R}{\mathbb R}
\newcommand{\dd}{\mathrm d}
\newcommand{\norm}[1]{\lVert #1\rVert}
\newcommand{\rg}{\rho_{\mathrm{root}}}
\setcounter{tocdepth}{1}
\DeclareMathOperator{\Span}{span}
\title{\bfseries 谱梯度法的端点动力学\\与有限记忆收敛}
\author{}
\date{2026 年 11 月 25 日}
\begin{document}
\maketitle
\begin{abstract}
同样沿负梯度走，为什么 Barzilai--Borwein（BB）方法的最坏渐近速率由活跃谱端点决定，而有限记忆最速下降（LMSD）的收敛分析需要整个梯度历史？本文从正定二次模型出发，推导平衡端点轨道、端点几何平均的最优权重、闭锥上的统一收缩，以及压缩 Ritz 扫掠的多项式表示。这些推导为阅读四篇公开预印本提供入口，并说明根渐近因子、统一指数估计和实际计算成本分别回答什么问题。
\end{abstract}
\smallskip
{\small\textbf{版本说明}\quad 根据四篇公开合著预印本独立编排的阅读稿。正文引用保留原作者和版本。整理与文字修订使用 AI 辅助，作者尚需逐项核阅。引用研究结论时请引用原始预印本。}
\tableofcontents

\section{一个二次模型和三种速率}

设 $H\in\R^{d\times d}$ 对称正定，$f(x)=\tfrac12 x^\top Hx-q^\top x$，$x_*=H^{-1}q$。记梯度 $g_k=Hx_k-q$。所有步长差异都体现在
\begin{equation}\label{e:gradient}
 g_{k+1}=(I-\alpha_kH)g_k.
\end{equation}
每个特征方向乘以 $1-\alpha_k\lambda$。因此，一个谱分量只要消失便不会重新出现。令 $a,b$ 为开始时实际出现的最小、最大特征值，记 $c=(b-a)/(b+a)$。

对一条非零初始梯度的轨道，$\rg=\limsup_{k\to\infty}(\norm{g_k}/\norm{g_0})^{1/k}$ 测量长期的每步衰减。固定某个 $\gamma$ 的\textbf{一致包络}则要求存在同一个 $C$，使类中所有允许的轨道和所有 $k$ 都满足 $\norm{g_k}\le C\gamma^k\norm{g_0}$。最后，实际成本还包括形成步长、作内积、向量正交化及处理有限精度等开销。迭代步数、梯度调用数和运行时间须各自记录。

例如，令 $0<c<1$。序列 $s_k=(k+1)c^k$ 的根因子是 $c$，在 $\gamma=c$ 时没有有限包络常数，在任意 $\gamma>c$ 时却有。另一个区别是初始化：一个任意大的正首步一般使 $\norm{I-\alpha_0H}$ 任意大，所以对任意首步的根渐近结论不能直接读成从 $k=0$ 开始、统一以 $\norm{g_0}$ 为基准的界。

\section{BB 轨道及其端点证书}

\subsection{延迟 Rayleigh 商与平衡轨道}

令 $s_{k-1}=x_k-x_{k-1}$、$y_{k-1}=g_k-g_{k-1}=Hs_{k-1}$。经典 BB1、BB2 的步长分别是
\[
 \frac{s^\top s}{s^\top y},\qquad
 \frac{s^\top y}{y^\top y}.
\]
当 $s=-\alpha_{k-1}g_{k-1}\ne0$，它们可以统一写为
\begin{equation}\label{e:rq}
 \tau_k=\alpha_k^{-1}
 =\frac{g_{k-1}^\top WHg_{k-1}}{g_{k-1}^\top Wg_{k-1}},
 \qquad W=I\text{ 或 }H.
\end{equation}
更一般地，固定正谱权重 $W=\omega(H)\succ0$ 也适用。$W^{1/2}$ 与 $H$ 交换，令 $\widetilde g=W^{1/2}g$，便把加权分式化为普通 Rayleigh 商，且 $\widetilde g$ 仍满足~\eqref{e:gradient}。两套范数在固定 $W$ 下等价，根因子不变。前因子换回 Euclidean 范数时可带来 $\sqrt{\kappa_2(W)}$。

在只含 $a,b$ 两端点分量的初始条件下，设两端点的加权能量相等，即
$\norm{W^{1/2}P_a g_0}^2=\norm{W^{1/2}P_b g_0}^2>0$。
匹配首步取 $\tau_0$ 为 $g_0$ 的同一个加权 Rayleigh 商，于是 $\tau_0=(a+b)/2$。这一步给出
\[
 P_ag_1=cP_ag_0,\qquad P_bg_1=-cP_bg_0 .
\]
两端平方同时乘 $c^2$，平衡保持。由~\eqref{e:rq} 递推，每一步都有 $\alpha_k=2/(a+b)$，且
\begin{equation}\label{e:orbit}
 g_k=c^k(P_ag_0+(-1)^kP_bg_0).
\end{equation}
因此 $\norm{g_k}/\norm{g_0}=c^k$ 恰好成立。平衡条件是在所用度量中写下的。特别是 BB2 一般不需要两端点的原始 Euclidean 分量相等。

\subsection{为什么一个几何平均能每步收缩}

当 $a<b$ 且两端点在首步后仍然存在时，后续 Rayleigh 商落在 $(a,b)$。记 $p=b/(a+b)$，定义
\[
 \mathcal L(g)=\norm{P_a g}^{2p}\norm{P_b g}^{2(1-p)},\qquad
 \Phi_p(\tau)
 =\Big(\frac{\tau-a}{\tau}\Big)^p
  \Big(\frac{b-\tau}{\tau}\Big)^{1-p}.
\]
沿~\eqref{e:gradient} 有 $\mathcal L(g_{k+1})=\Phi_p(\tau_k)^2\mathcal L(g_k)$。计算导数可得
\[
 \frac{\dd}{\dd\tau}\log\Phi_p(\tau)
 =\frac{ab-\tau((1-p)b+pa)}{\tau(\tau-a)(b-\tau)} .
\]
对于上述 $p$，分子只在 $(a+b)/2$ 过零，从正变负。两端极限为零，中点处 $\Phi_p=c$。令 $D(\tau)=\log c-\log\Phi_p(\tau)$，即得
\[
 \mathcal L(g_{k+1})=c^2\exp(-2D(\tau_k))\mathcal L(g_k),\qquad D\ge0 .
\]
中点给所有 $p\in(0,1)$ 同一个值 $c$。其他 $p$ 的导数在这里非零，因而在附近会取得严格大于 $c$ 的值。这证明了端点单项式中权重的唯一极小极大性。

这个证书测量端点的共同存活。如果某端点质量趋零，$\mathcal L$ 可能远小于 $\norm g^2$。处理全梯度时还要使用活跃谱和尾轨道论证。\href{https://arxiv.org/abs/2609.12084}{端点 Lyapunov 预印本}给出这些补充，并展示 Lyapunov 量下降而梯度范数不单调的轨道。

\subsection{尖锐上界还需要什么}

上面的显式轨道已经给出最坏因子的下界。对任意初值的上界要追踪活跃谱随时间的变化，尤其要处理质量靠近零的情形。公开的尖锐率预印本把延迟更新写为相邻归一化谱权重的动力学。在周期情形，选取仍有质量的端点并对对数伸缩因子求和，端点归一化比值在返回时抵消。对一般轨道，作者通过紧化和不变测度处理长期平均，再用半一致估计回到全轨道。由此得到有限维活跃区间的上界 $c$，并与~\eqref{e:orbit} 的可达性合并。

Hilbert 空间部分需要用标量谱测度代替有限谱的权重表。谱端点可以没有特征向量，因此精确的二端点轨道只在存在相应向量时直接适用。最优一致阈值的证明仍需处理靠近连续谱端点的初始数据。非线性部分将驻点附近的实际割线迭代与严格 Fréchet 线性化联系起来，控制 $\gamma>c$ 的局部包络，并从收敛尾轨道得到根因子界。这些扩展的完整假设和证明见\href{https://arxiv.org/abs/2608.07839}{尖锐率预印本}的正文及附录。

\section{从逐轨道收敛到统一指数界}

延迟意味着 $g_k$ 本身没有记住确定步长的全部信息。对 BB，可以在相容的连续梯度对上建一阶状态。记 $m,M$ 为 $H$ 的全谱最小与最大特征值。相容是说存在 $\alpha\in[1/M,1/m]$，使后一个梯度等于 $(I-\alpha H)$ 作用在前一个梯度的结果。允许任意正首步时，第一次通过 Rayleigh 商得到的受控更新会进入这个状态空间。若只知道真实运行会趋零，还需要验证锥中每个状态可由合法运行实现，并证明映射在终止边界连续。这是\href{https://arxiv.org/abs/2609.12100}{BB 闭锥预印本}中依赖算法的部分。

\begin{proposition}[齐次轨道的一致收缩，Shen--Hu 的单映射情形]
设 $\Omega$ 为有限维赋范空间中的闭锥，$T:\Omega\to\Omega$ 连续且正一次齐次。若对每个 $z$ 有 $T^kz\to0$，则存在 $C\ge1$ 和 $0<\gamma<1$，使 $\norm{T^kz}\le C\gamma^k\norm z$ 对一切 $z,k$ 成立。
\end{proposition}
\begin{proof}
若 $\Omega=\{0\}$，结论直接成立。其余情形，在紧集 $S=\{z\in\Omega:\norm z=1\}$ 上，对每个 $z$ 选正整数 $n(z)$ 使 $\norm{T^{n(z)}z}<1/2$。连续性给出邻域，在其中同一时刻仍满足严格不等式。取有限覆盖，并令 $N$ 为这些时刻的最大值。于是任一非零状态可在至多 $N$ 步内将范数缩小至少一半。

对中间状态重新选择时刻。设截至第 $k$ 步已经走完 $m$ 个这样的收缩段，未完成的余段少于 $N$ 步，则 $m\ge k/N-1$。紧性还给
$B=\max_{0\le j<N}\sup_{z\in S}\norm{T^jz}<\infty$。齐次性把球面估计乘回实际尺度，得到
$\norm{T^kz}\le B2^{-m}\norm z \le 2B(2^{-1/N})^k\norm z$。到达零时结论直接成立。
\end{proof}

这一论证给出的 $\gamma$ 未必最紧。它也依赖单位截面的紧性，不能直接移植到无限维。若重新选择的步号是整轮 LMSD 扫掠的端点，指数首先按扫掠计量，换为单次梯度更新还需要步数换算。

\section{LMSD 中的谱压缩和整轮多项式}

\subsection{同一历史可以编码多个步长}

设非终止运行上一轮的梯度历史矩阵为 $G$，秩为 $r\ge1$，令 $Q$ 为 $\operatorname{ran}G$ 的正交基矩阵。$B=Q^\top HQ$ 正定，Ritz 值是它的特征值 $\theta_1,\ldots,\theta_r$。下一轮依次取 $\alpha_j=\theta_j^{-1}$。因为 $H$ 固定且各因子都是它的多项式，从起点 $g$ 到终点的总作用是 $\Psi(H)g$，其中
\[
 \Psi(t)=\prod_{j=1}^{r}(1-t/\theta_j)
 =\frac{\det(Q^\top(H-tI)Q)}{\det(Q^\top HQ)}.
\]
这个公式只涉及 $\operatorname{ran}G$。历史矩阵出现确切秩退化时，$r$ 也应相应下降。使用数值容差截断是另一个算法选择，需要另外核对其有限精度稳定性。

\subsection{展成正权重组合}

在 $H$ 的一个正交特征坐标系内，令 $Q_I$ 是从 $Q$ 取下标集 $I$ 的 $r$ 行得到的方阵。对角矩阵与 Cauchy--Binet 公式给出
\[
 \det(Q^\top(H-tI)Q)=\sum_{|I|=r}(\det Q_I)^2\prod_{i\in I}(\lambda_i-t).
\]
令 $t=0$ 并归一化，得到
\begin{equation}\label{e:convex}
 \Psi(t)=\sum_{|I|=r}w_I
              \prod_{i\in I}(1-t/\lambda_i),\qquad
 w_I=\frac{(\det Q_I)^2\prod_{i\in I}\lambda_i}{\det(Q^\top HQ)},\quad
 \sum_I w_I=1 .
\end{equation}
所有 $w_I\ge0$。式~\eqref{e:convex} 把压缩矩阵的根与真实谱的可观测组合联系起来。因子本身在某些谱区间可以大于一；证明任务因此是控制反复生成的历史及其质量，而非要求每个扫掠对所有谱点都收缩。

有限终止也可从这里理解。在一个完整的梯度历史确实生成起点的活跃 Krylov 子空间时，若活跃谱点数不超过名义记忆阶数，限制算子的 Ritz 值就是该子空间上的所有活跃特征值。当前与下一轮梯度不会引入新的谱点。因而下一轮多项式包含消掉每个仍存谱点的因子。一般情况下，梯度可能一直有多于记忆阶数的活跃谱点。公开论文用行列式公式和谱归纳证明逐点收敛，再证明相容端点转移的连续性，最后应用闭锥上的统一收缩。\href{https://arxiv.org/abs/2609.12114}{加权 LMSD 预印本}给出了固定正谱权重下的共轭、终止分支以及从端点转为中间梯度估计所需的条件。

\section{怎样使用这些结论}

\begin{tabularx}{\linewidth}{@{}p{.22\linewidth}X@{}}
\toprule
阅读目的& 应核对的数学问题\\\midrule
判断速率是否尖锐 & 是否有可实现的见证轨道。区分任意正首步的单轨道结论与匹配首步的一致包络。\\
解释非单调轨道 & 看实际记忆状态、端点证书以及端点质量的消失。区分证书的大小、目标值和梯度范数。\\
分析有限记忆算法 & 对照名义记忆、代数秩、有效扫掠长度与更新时钟。记录证明需要哪个相容状态空间。\\
比较实际效率 & 连同步长估计、正交化、精度、停机规则和梯度调用一起计量。理论的根因子无法单独给出计时结论。\\
\bottomrule
\end{tabularx}

公开补充程序只用 Python 标准库，采用有理运算校验两条小例子。第一条分别对 BB1 与 BB2 验证匹配平衡轨道；第二条在一个三节点、二维压缩的例子中比较行列式恒等式的两边。程序是公式及实现的复核。一般定理的依据是上面的数学证明及所引论文。

\section*{参考文献}
\begin{itemize}
\item S. Yang and Y.-X. Yuan, \emph{The sharp worst-case asymptotic rate of the Barzilai--Borwein method in $\R^d$ and Hilbert spaces}, \href{https://arxiv.org/abs/2608.07839v1}{arXiv:2608.07839v1}, 2026.
\item S. Yang and Y.-X. Yuan, \emph{Lyapunov functions and $R$-linear convergence for quadratic Barzilai--Borwein dynamics}, \href{https://arxiv.org/abs/2609.12084v1}{arXiv:2609.12084v1}, 2026.
\item S. Yang and Y.-X. Yuan, \emph{$R$-linear convergence of Barzilai--Borwein methods on strictly convex quadratics via closed-cone homogeneous dynamics}, \href{https://arxiv.org/abs/2609.12100v1}{arXiv:2609.12100v1}, 2026.
\item S. Yang and Y.-X. Yuan, \emph{Global and unconditional $R$-linear convergence of rank-compressed weighted LMSD sweeps for strictly convex quadratics}, \href{https://arxiv.org/abs/2609.12114v1}{arXiv:2609.12114v1}, 2026.
\item J. Barzilai and J. M. Borwein, \emph{Two-point step size gradient methods}, IMA J. Numer. Anal., 8 (1988), 141--148. \href{https://doi.org/10.1093/imanum/8.1.141}{DOI}.
\item R. Fletcher, \emph{A limited memory steepest descent method}, Math. Program., 135 (2012), 413--436. \href{https://doi.org/10.1007/s10107-011-0479-6}{DOI}.
\item J. Shen and J. Hu, \emph{Stability of discrete-time switched homogeneous systems on cones and conewise homogeneous inclusions}, SIAM J. Control Optim., 50 (2012), 2216--2253. \href{https://doi.org/10.1137/110845215}{DOI}.
\item G. Ferrandi and M. E. Hochstenbach, \emph{Limited memory gradient methods for unconstrained optimization}, Numer. Algorithms, 99 (2025), 735--764. \href{https://doi.org/10.1007/s11075-024-01895-9}{DOI}.
\end{itemize}
\end{document}
